\documentclass[a4paper,12pt]{article}
\usepackage[english]{babel}
\usepackage{graphicx}
\usepackage{color}
\usepackage[separate-uncertainty = true,multi-part-units=single]{siunitx}
\sisetup{range-phrase=-,range-units=single}
\usepackage[margin=2cm]{geometry}
\usepackage{authblk}
\usepackage{amssymb}
\usepackage{amsmath}
\sloppy
\linespread{1}
\usepackage[numbers,sort&compress]{natbib}
\renewcommand\Affilfont{\fontsize{11}{15}\itshape}
\usepackage[labelfont=bf]{caption}

\begin{document}
	
	\title{Soap bubble lasers}
	\date{}
	\author[1,2]{Zala Poto\v cnik}
	\author[1,2,3]{Matja\v z Humar*}
	\affil[1]{Condensed Matter Department, J. Stefan Institute, Jamova 39, SI-1000 Ljubljana, Slovenia}
	\affil[2]{Faculty of Mathematics and Physics, University of Ljubljana, Jadranska 19, SI-1000, Ljubljana, Slovenia}
	\affil[3]{CENN Nanocenter, Jamova 39, SI-1000 Ljubljana, Slovenia}
	\affil[ ]{*E-mail: matjaz.humar@ijs.si}
	\affil[ ]{Keywords: soap bubble, whispering gallery modes, laser, liquid crystals, electric field sensor, pressure sensor}
	\maketitle
	

	\begin{abstract}
		
		Soap bubbles are simple, yet very unique and marvelous objects. They exhibit a number of interesting properties such as beautiful interference colors and the formation of minimal surfaces. Various optical phenomena have been studied in soap films and bubbles, but so far they were not employed as optical cavities. Here we demonstrate, that dye doped soap or smectic liquid crystal bubbles can support whispering gallery mode lasing, which is observed in the spectrum as hundreds of regularly spaced peaks, resembling a frequency comb. The lasing enabled the measurement of size changes as small as \SI{10}{nm} in a millimeter-sized, \SI{\sim 100}{nm} thick bubble. Bubble lasers were used as extremely sensitive electric field sensors with a sensitivity of \SI{11}{Vm^{-1}Hz^{-1/2}}. They also enable the measurement of pressures up to a \SI{100}{bar} with a resolution of \SI{1.5}{Pa}, resulting in a dynamic range of almost $10^7$. By connecting the bubble to a reservoir of air, almost arbitrarily low pressure changes can be measured while maintaining an outstanding dynamic range. The demonstrated soap bubble lasers are a very unique type of microcavities which are one of the best electric field and pressure microsensors to date and could in future also be employed to study thin films and cavity optomechanics.
		
	\end{abstract}
	
	
	A soap bubble is made of a thin film composed of water and surfactants, which encloses air and forms a spherical shape \cite{isenberg1978science}. Bubbles are interesting from a variety of perspectives including mathematics, physics, chemistry and even biology, due to their similarity to the biological membranes. Soap bubbles were studied in terms of the interference colors \cite{afanasyev2011measuring}, geometry \cite{almgren1976geometry,taylor1976structure}, fluid motion \cite{debregeas1998life,meuel2013intensity}, mechanical oscillations \cite{kornek2010oscillations} and recently, branched flow of light \cite{patsyk2020observation}. A bubble can also be made purely from surfactant-like molecules. Specifically, smectic liquid crystals, which molecules form well defined molecular layers, are used for this purpose \cite{stannarius1998self,may2012dynamics,clark2017realization,lopez2011nematic}. The mixture of soap and water can actually also form liquid crystalline (lyotropic) phases. In fact, the origin of the word "smectic" is related to soap. Smectic bubbles have some very unique properties. Since the smectic liquid crystal molecules form well defined layers, the film thickness is quantized, that is to say, the film has an integer number of molecular layers. For example, bubbles with a completely uniform thickness of \SI{11}{nm} as large as \SI{1}{cm} in diameter were made \cite{stannarius1998self}. Therefore, quite uniquely, the ratio of the bubble size to its thickness can be as large as 6 orders of magnitude. The layered structure of the smectic bubbles makes their thickness completely stable and enables virtually infinite bubble lifetime as long as the air volume inside the bubble is kept constant. 
	
	Optical resonances called whispering gallery modes (WGMs) are formed when the light is trapped in a spherical object due to multiple total internal reflections and circulates near the surface of the sphere. WGMs were studied in various geometries including solid hollow cavities in the form of glass microbubbles \cite{sumetsky2010optical,watkins2011single} and glass capillaries \cite{shopova2007optofluidic}, which were employed for for sensing applications \cite{ward2014hollow} and lasing \cite{shopova2007optofluidic,chen2016lasing}. However, WGMs were not studied until now in soap bubbles.
	
	
	Here we show that dye doped soap and smectic bubbles can support WGM lasing. Due to their fluid nature the bubbles are very soft compared to their glass counterparts, which influences the lasing and enables some unique applications.
	
	
	\begin{figure}[htbp]
		\begin{center}
			\includegraphics[width=11cm]{fig1.pdf}
			\caption{\textbf{A soap bubble formed at the end of a capillary.} (a) Scheme of the experimental configuration. A dye doped soap bubble is inflated at the end of a horizontal capillary and illuminated by a laser from below. The soap film is composed of a layer of water, surfactant molecules and fluorescent dye molecules. (b) A soap bubble in reflected light. Interference colors are visible. (c) Fluorescence image of a dye doped bubble.}
			\label{fig1}
		\end{center}
	\end{figure}
	
	
	\subsection*{Lasing of regular soap bubbles}
	
	Millimeter-sized soap bubbles doped with a fluorescent dye, inflated at the end of a capillary and illuminated with a pulsed laser were used to demonstrate the laser emission (Fig. \ref{fig1}a). To create the bubbles, a capillary was dipped into the soap solution and the air pressure in the capillary was briefly increased to inflate the bubble (Supplementary Video 1). The soap bubble could be later inflated or deflated to reach the desired size (\SIrange{0.4}{4}{mm}). Based on the interference colors \cite{afanasyev2011measuring} observed in reflection (Fig. \ref{fig1}b) the thickness of the soap film was typically in the range \SIrange{100}{800}{nm} and changed slowly in time as the bubble aged. The soap bubbles had a relatively uniform fluorescence intensity, except at the end of the capillary where there was a larger amount of the solution (Fig. \ref{fig1}c). 
	
	
	
	\begin{figure}[htbp]
		\begin{center}
			\includegraphics[width=16cm]{fig2.pdf}
			\caption{\textbf{Soap bubbles emitting laser light.} (a) A soap bubble illuminated in the center (red cross) by a pump laser. The drawing shows the gain regions (red) which are formed at the positions where the laser beam (green) passes through the soap film. The WGMs circulating in any vertical plane (black lines) pass through the two gain regions and are therefore excited. (b) Another bubble but illuminated at its rim, which generates a bright ring of circulating light. The gain region in this case is in the shape of a vertical patch. WGMs circulating in a narrow band in a vertical plane experience the most gain and are therefore observed as a bright ring. (c) A typical emission spectrum when a bubble is pumped above the lasing threshold. (d) Lasing intensity summed over a \SI{5}{nm} wide spectrum range as the input laser pulse energy was increased shows a typical threshold behavior. The dashed lines are a guide to the eye. (e) Spectra of consecutive pump laser pulses. Displaying a smaller wavelength range reveals a shift of the modes towards longer wavelengths. (f) Experimental setup for the free floating soap bubbles. (g) A photo of a larger floating soap bubble, which is illuminated by the pump laser, and emits laser light. (h) Spectrum from a \SI{\sim2}{cm} diameter free floating bubble.}
			\label{fig2}
		\end{center}
	\end{figure}
	
	
	When a bubble was illuminated by a pulsed laser it emitted laser light due to the circulation of WGMs. When the center of the bubble was pumped, laser light was generated in all vertical planes except the ones intersecting with the capillary. This was observed as a bright rim, except on the opposite side of the capillary (Fig. \ref{fig2}a). When the bubble was illuminated at an edge, the generated light preferentially circulated in one plane, which was observed as a bright narrow ring (Fig. \ref{fig2}b).
	
	
	Sharp spectral lines were present in the spectrum of the emitted light (Fig. \ref{fig2}c). A clear lasing threshold was observed when the pump laser energy was increased (Fig. \ref{fig2}d). Typical thresholds were in the order of a few \SI{}{\micro J}.  The WGMs can be approximated as $2\pi r n_{\mathrm{eff}}= \lambda l$, where $r$ is the radius of the bubble, $n_{\mathrm{eff}}$ is the effective refractive index of the mode, $\lambda$ is the wavelength of the mode and $l$ is the azimuthal mode number. Since the bubbles were usually relatively large (\SI{\sim1}{mm}) compared to the wavelength of the light, as a first approximation, the light propagation can be described as a flat slab waveguide \cite{kawano2004introduction}. For a typical bubble with a thickness of \SI{600}{nm} the effective refractive indices were 1.319 and 1.307 for first radial order TE and TM modes, respectively. For a typical soap bubble thickness not only the fundamental, but also higher guided modes may be present. In the lasing spectrum no regularly spaced peaks could be observed. This is probably due to a large number of possible modes and irregularities of the soap film. With each pulse of the pump laser the spectrum changed slightly (Fig. \ref{fig2}e). Some modes disappeared, while some others appeared, but there was a general trend towards a shift towards longer wavelengths. The shifting of the peaks in time was probably caused by a combination of size, thickness and refractive index changes. However, from the spectrum alone we can not determine how each of the three parameters contributes to the shift in the spectrum.
	
	Lasing was also achieved in free floating soap bubbles. A bubble was blown into a tank filled with CO$_2$ so that it was floating on the surface of the CO$_2$. The bubble was illuminated by the pump laser by using an optical fiber and a lens (Fig. \ref{fig2}f and g). The same lens and optical fiber were also used to collect the generated light and send it to the spectrometer. Spectra with sharp spectral lines corresponding to lasing were also observed in this case (Fig. \ref{fig2}h).
	
	
	\subsection*{Lasing of smectic bubbles}
	
	To better control the thickness and the refractive index of the film, bubbles made out of smectic liquid crystal were employed. Specifically, 8CB was used, which has a smectic A phase at room temperature. Under transmitted light the smectic bubbles looked completely uniform (Fig. \ref{fig3}a) indicating that the whole surface had thickness of equal number of molecular layers (Fig. \ref{fig3}b). If some region had one more or one less molecular layer this would result in a well visible island or hole, respectively \cite{clark2017realization}. The thickness of the bubbles could be approximately controlled by how fast they were inflated, with the faster inflation resulting in a thinner bubble. The thickness was estimated by the amount of transmitted light \cite{stannarius1999self}. Imaging a bubble in between crossed polarizers \cite{mirri2015situ} indicates an uniform radial orientation of the molecules (Fig. \ref{fig3}c and d).  
	
	\begin{figure}[htbp]
		\begin{center}
			\includegraphics[width=17cm]{fig3.pdf}
			\caption{\textbf{Lasing of smectic bubbles and islands.} (a) A smectic bubble in transmitted light. (b) Scheme of the molecular structure of the film. Drawn is an example of a three layer smectic film composed of ordered liquid crystal and dye molecules. (c) Same bubble under crossed polarizers. The typical bright cross is observed indicating uniform molecular orientation either parallel or perpendicular to the bubble surface. (d) Same bubble with an additional waveplate inserted between the polarizers. From the resulting colors it can be deduced that the molecules are oriented perpendicular to the surface, that is in the bubble radial direction. (e) Image of a lasing smectic bubble. The cross indicates the pump laser beam position. (f) Frequency comb-like spectrum of a lasing smectic bubble \SI{1.75}{mm} in diameter. (g) Shematics of an island floating on a smectic bubble supporting WGM lasing. (h) Two smectic islands (indicated by the arrows) on a bubble. (i) Image of a lasing island observed as a red ring (white arrow) on the surface of the bubble (indicated by the red arrows). (j) Lasing spectra of two smectic islands of a different diameter $d$. 
			}
			\label{fig3}
		\end{center}
	\end{figure}
	
	 When pumping with a pulsed laser, the dye doped smectic bubbles emitted laser light visible as a ring and two bright spots (Fig. \ref{fig3}e). The smectic bubbles were extremely stable and the experiments on a single smectic bubble could be performed for up to \SI{30}{min}. The emission spectrum above the lasing threshold was very regular with equally spaced spectral lines, resembling a frequency comb (Fig. \ref{fig3}f). The emission typically spanned \SIrange{10}{20}{nm}. For a \SI{1.8}{mm} diameter bubble the average free spectral range (FSR) was \SI{0.06}{mm} and the spectrum contained $\sim 250$ lasing peaks. The contrast between the fluorescent background and the lasing peaks in the spectrum was excellent with the background being practically zero. By measuring the emission intensity through a polarizer the modes were identified to have TE polarization.
	
	
	From the measured FSR and the diameter of the bubble determined from the image (\SI{1.90 \pm 0.02}{mm}), the effective refractive index is calculated to be $1.06 \pm 0.01$. For TE modes, using equations for a slab waveguide, this corresponds to a bubble thickness of \SI{57 \pm 5}{nm}. Typically the bubbles employed in this study had a thickness in the range of \SIrange{30}{120}{nm}, which matches the thickness determined from the intensity of the transmitted light. Thinner smectic bubbles did not generate laser emission due to too low effective refractive index causing significant radiative light leakage and the fact that a thin layer also results in very few dye molecules to be present, resulting in a low gain. At a typical thickness of \SI{100}{nm} the TM modes have a lower effective refractive index (1.07) than TE modes (1.15) (Supplementary Fig. 4), which is also the reason they have a higher radiative loss and are therefore not observed in the lasing spectrum. At this typical bubble thicknesses higher radial modes are not allowed. This is a significant advantage over bulk WGM microcavities which will typically display lasing of a number of different modes when their size is large. On the contrary, the lasing spectrum of smectic bubbles is very regular, frequency comb-like. Further, the size of smectic bubbles can be changed in real time, within a few seconds, by inflating or deflating them in at least the size range \SIrange{0.5}{5}{mm} resulting in a large tunability of the FRS in the range of \SIrange{0.02}{0.2}{nm}.
	
	
	When the volume of the bubble is increased or decreased very quickly, stable circular floating islands with larger thickness than the rest of the bubble can be formed (Fig \ref{fig3}g and h) \cite{clark2017realization,stannarius1998self}. The islands can also be created by illuminating a thicker bubble ($\gtrsim \SI{200}{nm}$) with a pulsed laser, where each pulse can create one island. The size of the islands can be from a few micrometers up to being as large as the bubble. Islands larger than \SI{\sim 100}{\micro m} emitted laser light upon excitation with a pulsed laser by the circulation of light around the perimeter of the island giving rise to WGM lasing (Fig. \ref{fig3}i), with periodic peaks in the spectrum (Fig. \ref{fig3}j). A typical effective refractive index calculated for the island size and FSR was $1.3 \pm 0.1$. This enabled stable laser emission with periodic lasing peaks with a wide range of FSR values (\SIrange{0.1}{0.9}{nm}) depending on the size of the islands (Fig. \ref{fig3}j).
	
	
	
	\subsection*{Smectic bubble lasers as electric field and pressure sensors}
	
	The smectic bubble lasers are very soft and therefore sensitive to external factors, which can change their size and shape. Simultaneously these changes can be measured very precisely via the lasing spectrum. Therefore, this makes soap bubble lasers excellent sensors.
	
	
	The changes in the bubble size were measured by following shifts of the lasing peaks. The relative size change is approximately equal to the relative wavelength shift: $\Delta d/d \approx \Delta \lambda/\lambda$. From the size measurement in time obtained in this way it is visible that the size of the bubble slowly decreases in time (Fig. \ref{fig4}a) due to the diffusion of air molecules through the thin wall \cite{ishii2009gas}. For real applications the size decrease could be compensated by pumping additional air into the bubble, stabilizing the size, or simply reinflating the bubble when it becomes too small, thus enabling almost indefinitely long measurements. The size change rate was \SI{1.1}{\micro m/s} for a \SI{1.9}{mm} bubble. The time dependence was fitted to the equation for the air diffusion (Supplementary Section 1) and subtracted from the data. The remaining residuals are extremely small with a standard deviation of only \SI{10}{nm} (Fig \ref{fig4}b). This means that size changes as small a \SI{10}{nm} can be measured, corresponding to an exceptional relative accuracy of $5 \times 10^{-6}$. A part of the observed noise may come from the environment, such as air pressure fluctuations and air currents, so that the actual optically limited measurement noise may be even lower than observed. To be able to track the spectral shifts, the shift in between two measurements (laser pump pulses) should be smaller than half the FSR. In the cases the changes were faster than this, FSR was measured instead, which enabled the measurement of arbitrarily large and fast size changes with a still excellent accuracy of \SI{50}{nm}. By taking into account the typical effective refractive index, the absolute size of the bubble can be measured via the FSR value. 
	
	
	\begin{figure}[htbp]
		\begin{center}
			\includegraphics[width=17cm]{fig4.pdf}
			\caption{\textbf{Electric field and pressure measurement with smectic bubble lasers.} (a) Change in the diameter of a bubble in time measured by following the lasing peaks. The initial size was measured from the FSR. The bubble size was slowly decreasing due to air diffusion through the thin wall. Since the size measurement is very precise, the plot looks like a perfectly smooth line. (b) When the decrease in size is subtracted from the data, only an extremely small noise remains. (c) Shift of the lasing modes when an electric field is increased continuously to \SI{3}{V/mm} and back to zero. The continuous decrease in the size due to diffusion was subtracted from the data. (d) Change of the bubble size measured from the lasing spectrum as the pressure outside the bubble was increased by \SI{400}{Pa} above the atmospheric pressure. $V_c/V_b+1\approx20$ for this particular case. (e) The calculated pressure sensitivity of a \SI{2}{mm} bubble as a function of total volume relative to the volume of the bubble ($V_c/V_b+1$). The dotted line would be for a bubble with zero surface tension. The red circle represents the maximum $V_c$ used in the experiments. (f) Maximum measurable positive and negative pressures. (g) The dynamic range for positive pressure, that is the ratio between the highest and the lowest measurable pressures.}
			\label{fig4}
		\end{center}
	\end{figure}
	
	
	
	The bubble lasers were employed as electric field sensors and electrically tunable laser sources. Already applying less than \SI{0.1}{V} across the bubble deformed it \cite{ishii2011anomalous} and caused a measurable shift in the laser peaks (Fig. \ref{fig4}c). The spectral shift was \SI{0.08}{nm(V/mm)^{-1}}, which is about 10 times larger than best previous reports using WGM microcavities \cite{kiraz2008spectral,ali2012high,humar2009electrically}. The minimum measurable field determined from the measurement noise and the sensitivity was \SI{11}{Vm^{-1}Hz^{-1/2}}. With this, our sensors outperforms practically all micro DC electric field sensors including sensors based on microelectromechanical systems \cite{kainz2018distortion,han2022micro}, piezoelectric \cite{xue2019piezoelectric}, optical \cite{zhang2021review,calero2019ultra,zhu2015high} and nitrogen-vacancy centers in diamond \cite{dolde2011electric,chen2017high}. Further, in contrast to some other electric field sensors, the bubbles are not conducting and therefore do not distort the measured electric field.
	
	
	The compressibility of the air within the bubble and the ability to measure the size extremely precisely enables the use of the bubbles as pressure sensors. The spectral shift for an isothermal pressure change is $\Delta \lambda/\lambda \approx - \Delta p/3 p_0$, where $\Delta p$ is the pressure change and $p_0$ is the initial pressure. In all the following measurements and calculations the initial pressure was equal to the atmospheric pressure. This gives a size independent sensitivity $\Delta \lambda/\Delta p$ of \SI{2e-3}{nm/Pa}. This is several orders of magnitude more than previously reported for WGM microcavities, such as hollow polymer spheres (\SI{2e-6}{nm/Pa}) \cite{ioppolo2007pressure} and a glass microbubbles (\SI{4e-7}{nm/Pa}) \cite{henze2011tuning}. 
	
	To test the performance of the pressure sensor, the pressure around the bubble was controllably changed and the diameter of the bubble was calculated from the measured FSR value. When the pressure was in increased by \SI{400}{Pa} above the atmospheric pressure the bubble diameter decreased by \SI{52}{\micro m} (Fig. \ref{fig4}d). In all experiments the bubble was connected to a capillary which provided an additional fixed volume of air. Since a larger volume of air is easier to compress, this increases the pressure sensitivity (Supplementary Section 2). By having a very large additional volume, an arbitrarily high sensitivity can be in principle achieved (Fig. \ref{fig4}e). At small additional volumes the sensitivity is inversely proportional to the total volume of the system  $V_c+V_b$, or more conveniently written as the total volume relative to the volume of the bubble $V_c/V_b+1$, where $V_c$ is the additional volume in the capillary and $V_b$ is the initial volume of the bubble. In the experiments up to $V_c/V_b+1=60$ could be easily achieved providing a sensitivity of \SI{0.12}{nm/Pa}. Above approximately $V_c/V_b+1=1000$ the surface tension starts to have a significant contribution. Above some critical additional volume, $V_c^{crit}/V_b+1=3130$ for this case, the sensitivity diverges, since the bubble becomes unstable and collapses by itself due to the Laplace pressure. The bubble also collapses at any $V_c>0$ if a large enough positive outside pressure is applied. This limits the maximum measurable pressure (Fig. \ref{fig4}f). The maximum measurable positive pressure is also limited by the minimum size of the bubble which is still capable of lasing (\SI{\sim 0.5}{mm}). The maximum measurable negative pressure is limited by the maximum size of the bubble for which the spectral lines can still be distinguished by the spectrometer (\SI{\sim 5}{mm}). 
	
	
	The smallest detectable pressure changes (Supplementary Fig. 9) were calculated from the sensitivity and the smallest measurable wavelength shift. By increasing $V_c$ the smallest detectable pressure change can be almost arbitrarily small, however, at the expense of the maximum measurable pressure.  For a stand alone bubble ($V_c=0$) the smallest measurable pressure is \SI{1.5}{Pa}, while the largest positive pressure is almost \SI{e7}{Pa}, that is \SI{100}{bar}. Namely, a pressure of \SI{e7}{Pa} would shrink a bubble from \SI{3}{mm} to \SI{0.65}{mm}, well with the lasing regime. This gives an exceptionally large dynamic range of almost $10^7$. For the largest additional volume used in the experiments ($V_c/V_b+1=60$) the smallest measurable pressure is \SI{0.024}{Pa}, while the largest positive and negative measurable pressures are \SI{1400}{Pa} and \SI{-1.9e4}{Pa}, respectively. Therefore by only changing the additional volume $V_c$, the pressure measurement range can be tuned in a huge range while preserving outstanding dynamic range of $\sim10^5$ (Fig. \ref{fig4}g).
	
	
	These values far exceed the performance of the best pressure sensors of a comparable size \cite{sahota2020fiber,vorathin2020review,domingues2018cost,zhu2013graphene,wang2016graphene,javed2019review,lin2018ultra,vsivskins2020sensitive,chen2019ultrahigh}. Especially the dynamic range of the bubble lasers is typically several orders of magnitude larger than other pressure microsensors demonstrated till now. Further, the significant advantage of the bubbles over pressure sensors using solid transducers, such as strain gauges and diaphragms, is that the smectic layer is very soft as well as practically infinitely stretchable, due to its fluid nature. The bubbles also do not experience any material fatigue. Further, the bubble pressure sensors do not need any calibration, since their sensitivity depends only on the compressibility of the air (for $V_c=0$). This is in stark contrast to almost all sensors, where variability in the mechanical properties and nonlinearities of the material (e.g. diaphragms), hysteresis and manufacturing tolerances play a crucial role in the accuracy of these sensors. 
	
	
	
	\subsection*{Conclusion and outlook}
	
	In conclusion, we demonstrated lasing in soap and smectic bubbles. Compared to solid or droplet microcavities, bubbles are made of a very thin film of fluid, resulting in truly unique optical and mechanical properties. The size of a bubble laser can easily be changed by a factor of 10 or more in real time. Due to free space light coupling, and their relatively large size, the excitation and collection of light is extremely easy and convenient. Only a multimode fiber and a simple lens are needed. The smectic bubbles have a completely uniform thickness down to a molecular level. The thin walls support only a single optical mode, therefore they can be employed as tunable laser sources with frequency-comb like emission, with orders of magnitude larger tunability compared to solid microlasers. Small thickness also results in small weight and therefore almost perfectly spherical shape. The bubbles are extremely stable in time and could in principle survive indefinitely. Lastly, the bubbles are very easy to form, not requiring any complicated manufacturing procedures and if a bubble is destroyed, a new one can be formed within seconds. These outstanding properties enable a number of prominent applications. Their exceptional sensitivity enables sensing of electric field and pressure with record high sensitivity, resolution and dynamic range. In future, they could be also used to measure other quantities, which change the bubble shape, such as air flow or magnetic field by doping with magnetic particles. Apart from the demonstrated applications, the bubble lasers could be used in future to study thin films and basic phenomena such as for example cavity optomechanics \cite{aspelmeyer2014cavity,kim2013cavity,asano2016observation}.
	
	
	\subsection*{Materials and Methods}
	
	A mixture of water, glycerol and liquid hand soap (containing sodium laureth sulfate) in a 2:1:1 volume ratio was used to make the soap bubbles ($n=1.364$). Alternatively, 2.5\% of sodium dodecyl sulfate dissolved in 1:1 volume ratio of water and glycerol was used. For the gain, 0.1\% of fluorescein sodium salt was dissolved in the above mixtures. For the smectic bubbles, 4'-octyl-4-biphenylcarbonitrile (8CB) liquid crystal doped with 0.2\% Pyrromethene 597 (Exiton) was used. The bubbles were inflated via a thin tube connected to a glass syringe (Hamilton, \SI{1.0}{ml}). The syringe was mounted onto a microfluidic syringe pump (New Era Pump Systems, NE-1002X). For electric field tuning the bubble was placed inbetween two flat electrodes with an area of $18 \times$\SI{20}{mm} and a spacing of \SI{5.46}{mm}. One of the electrodes had a hole with a diameter of \SI{2.6}{mm} through which the pipette tip was inserted so that the bubble was in the center between the electrodes. For pressure measurements the pipette tip was inserted through a hole into a transparent container, which was connected to a pressure controller (Elveflow, OB1, \SIrange{0}{20}{kPa}). The bubbles were observed with an inverted microscope (Nikon Ti2) through a $4\times$, \SI{0.13}{NA} objective. The bubbles were pumped with a nanosecond pulsed optical parametric oscillator (Opotek, Opolette 355) at \SI{494}{nm} for the soup bubbles and \SI{525}{nm} for the smectic bubbles, at a repetition rate of \SI{20}{Hz}. The resulting fluorescent light was captured by a high resolution spectrometer (Andor Shamrock SR-500i) at \SI{0.007}{nm} spectral resolution.
	
	
	\subsection*{Author Contributions}
	Z.P. performed the experiments. Z.P. and M.H. analyzed the data and wrote the manuscript. M.H. designed and supervised the study.
	
	\subsection*{Acknowledgments}
	This project has received funding from the European Research Council (ERC) under the European Union’s Horizon 2020 research and innovation programme (grant agreement No. 851143) and from Slovenian Research Agency (ARRS) (J1-1697 and P1-0099).
	
	\subsection*{Competing interests}
	The authors declare no competing interests.
	
	\subsection*{Data availability}
	All data that support the plots within this paper and other findings of this study are available from the corresponding author upon reasonable request.
	
	\subsection*{Code availability}
	The codes that support the findings of this study are available from the corresponding author upon reasonable request.
	
	
	
\clearpage

\setcounter{figure}{0}

\section*{Supplementary Information}


\begin{figure}[h]
	\begin{center}
		\renewcommand\figurename{Supplementary Figure}
		\includegraphics[width=10cm]{FigS1.pdf}
		\caption{A soap bubble pumped at the center. Light only circulates in planes intersecting with the pump spot, but not intersecting the capillary, which can be seen as a dark patch on the opposite side of the capillary (red arrow).}
		\label{fig1}
	\end{center}
\end{figure}

\clearpage

\begin{figure}[h]
	\begin{center}
		\renewcommand\figurename{Supplementary Figure}
		\includegraphics[width=13cm]{FigS2.pdf}
		\caption{A dye doped soap bubble at the end of a capillary. (a) In white light interference colors can be observed. (b) The same bubble when illuminated by the pump laser. The arrow shows the direction of the incoming laser beam. A bright ring, which crosses the laser spot, is formed due to circulating light.}
		\label{fig2}
	\end{center}
\end{figure}

\clearpage

\begin{figure}[h]
	\begin{center}
		\renewcommand\figurename{Supplementary Figure}
		\includegraphics[width=10cm]{FigS3.pdf}
		\caption{Spectrum emitted by a soap bubble attached to the end of a capillary when the pulse energy of the pump laser is increased. At lower energies only fluorescence is observed, while at higher energies sharp peaks appear.}
		\label{fig3}
	\end{center}
\end{figure}


\clearpage

\begin{figure}[h]
	\begin{center}
		\renewcommand\figurename{Supplementary Figure}
		\includegraphics[width=16cm]{FigS4.pdf}
		\caption{Effective refractive indices for different modes calculated for a slab waveguide. The shaded areas in both panels correspond to the approximate thickness range of the bubbles used in the experiments. Calculating the positions of WGMs by using the effective index for flat slab is accurate to a few percent compared to the WGMs simulations \cite{zhang2017chip}. (a) Modes for a soap film at the wavelength of \SI{555}{nm}. For a typical bubble thickness there are a number of modes allowed. (b) Modes for a smectic film at the wavelength of \SI{610}{nm}. Since, the smectic film is birefringent ($n_\mathrm{o}=1.51$ and $n_\mathrm{e}=1.68$), the bulk refractive indices of TE and TM modes are different (dashed lines). For a typical smectic bubble thickness (\SIrange{30}{120}{nm}) only TM$_0$ and TE$_0$ exist. However, TE$_0$ has a significantly larger effective refractive index than TM$_0$ and is the only mode lasing.}
		\label{fig4}
	\end{center}
\end{figure}


\clearpage

\begin{figure}[h]
	\begin{center}
		\renewcommand\figurename{Supplementary Figure}
		\includegraphics[width=15cm]{FigS5.pdf}
		\caption{Smectic island lasers. (a) Two examples of smectic islands (regions of higher thickness) floating on a smectic bubble. Each island has a different thickness, therefore each displays a different interference color. The smectic bubble (marked by the red arrows) is very thin and despite the enhanced contrast of the image, it is barely visible. (b) The islands can also be created by illuminating a thicker bubble ($>\SI{200}{nm}$) with a pulsed laser (black cross), where each pulse can create one island. The islands move away from the laser spot (green arrow). (c) Two examples of lasing from smectic islands, which are visible as bright rings.}
		\label{fig5}
	\end{center}
\end{figure}

\clearpage


\begin{figure}[h]
	\begin{center}
		\renewcommand\figurename{Supplementary Figure}
		\includegraphics[width=10cm]{FigS6.pdf}
		\caption{Lasing spectrum from a \SI{1.75}{mm} smectic bubble in time. The color represents the intensity. The white line highlights the shifting of a single lasing mode in time due to the decreasing size of the bubble.}
		\label{fig6}
	\end{center}
\end{figure}


\clearpage

\begin{figure}[h]
	\begin{center}
		\renewcommand\figurename{Supplementary Figure}
		\includegraphics[width=10cm]{FigS7.pdf}
		\caption{Data from the same bubble as in main text Fig. 4a and b. In this case the size of the smectic bubble was calculated from the FSR. The resulting decrease in the size of the bubble was fitted to Eq. \eqref{R(t)_approx} and subtracted from the data. The remaining residuals have a standard deviation of \SI{50}{nm}.}
		\label{fig7}
	\end{center}
\end{figure}


\clearpage

\begin{figure}[h]
	\begin{center}
		\renewcommand\figurename{Supplementary Figure}
		\includegraphics[width=10cm]{FigS8.pdf}
		\caption{Diameter of a bubble in time measured from FSR. The pressure outside the bubble is increased to \SI{400}{Pa} above the atmospheric pressure and then decreased to \SI{240}{Pa}. The continous decrease in size is due to air permeation through the smectic layer.}
		\label{fig8}
	\end{center}
\end{figure}


\clearpage	

\section{Air permeation through smectic bubbles}
The permeation of gas through the smectic film is proportional to the pressure difference between the sides of the film and the surface of the film. For a complete sphere, the gas volume loss
\begin{equation}
	\dot{V}=4\pi R^2 \dot{R}
\end{equation}
is proportional to the surface area, where $R$ is the radius of the bubble. In an approximation of a small pressure difference between the sides of the film, we assume a linear dependence of the volume loss on Laplace pressure $p_L=\frac{4 \sigma}{R}$, where $\sigma$ is the surface tension. Then
\begin{equation}
	4\pi R^2 \dot{R} \approx -\gamma 4\pi R^2 \frac{4 \sigma}{R},
\end{equation}
where $\gamma(d,T)$ is a phenomenological parameter of gas permeation describing the volume of gas passing through the film per pressure difference and area, where $d$ is the thickness of the smectic film and $T$ is the temperature. The solution of this differential equation is
\begin{equation}
	R^2(t) = R_0^2 - gt, \label{R(t)_approx}
\end{equation}
with the initial radius $R_0=R(0)$ and the fitting parameter $g(d,T)=8\sigma\gamma $. This equation was used to the data in main text Fig. 4a and c. This dependence holds as long as the bubble is large compared to the radius of the capillary opening ($r_c$). The first-order correction takes into account  that the area of the bubble is smaller due to the capillary opening. The subtracted area is approximated by a disc of radius $r_c$.
This results in a corrected equation for the size time dependence \cite{li2002gas}:
\begin{equation}
	R^2+\frac{r_c^2}{4}\ln{\left(R^2-\frac{r_c^2}{4}\right)}=R_0^2+\frac{r_c^2}{4}\ln{\left(R_0^2-\frac{r_c^2}{4}\right)}-gt.
	\label{R(t)}
\end{equation}

The experimental data of the smectic bubble radius shrinkage over time were fitted with Eq. \eqref{R(t)}.  The fitted parameters are $R_0=$ \SI{592}{\um}, $r_c=$ \SI{336}{\um} and $g=3.85\cdot 10^{-10} \mathrm{m^2/s}$. The gas permeation parameter $\gamma$ obtained from the fit is $\gamma=(2.0\pm0.1)\cdot 10^{-9} \mathrm{m/Pa\,s}$, where the surface tension value for 8CB smectic is assumed to be $\sigma=$ \SI{0.024}{N/m} \cite{stannarius1997surface}.
The good agreement of data and fit indicates that there was no leakage of air through the apparatus for bubble inflation, since the time dependence of the radius would be different in the case of leakage ($R^4 - R_0^4 \propto t$) \cite{li2002gas}.

The permeation of gas through smectic film depends strongly on the film thickness and temperature, as well as on the species of the gas molecules \cite{ishii2009gas, li2002gas}.
The permeation of gas molecules follows Fick's law $J=D\nabla C$, which describes the flow of a gaseous solute within a liquid solvent, where $J$ is the flow rate, $D$ is the diffusion constant, and $C$ is the space-dependent solute concentration.
The gas concentration in the surface layer of the bubble in contact with the gas can be expressed as $C=Sp$ using gas solubility in the smectic and pressure.
Combining this with Fick's law and considering the one-dimensional geometry of the flow, we get:
\begin{equation}
	J=P\,\frac{\Delta p}{d},
\end{equation}
where $P=DS$ is the permeability coefficient of the gas through the smectic film and depends on the diffusion constant and the solubility of the gas molecules in the smectic, $\Delta p$ is the pressure difference between the sides of the film, and $d$ is the film thickness \cite{ishii2009gas}.

If we combine the measured value of $\gamma$ with the permeability coefficient $P$ from the model, we can write:
\begin{equation}
	\gamma=P/d.
\end{equation}
Using the known permeability coefficients $P$ for different gases from the study \cite{ishii2009gas}, we can calculate the thickness of the observed bubble.
Assuming that air is composed of 78\% nitrogen and 22\% oxygen, we use the permeability coefficients $P_{\mathrm{N_2}}=(3.8\pm0.3)\cdot 10^{-15}\,\mathrm{mol\,Pa^{-1}m^{-1}s^{-1}}$ and $P_{\mathrm{O_2}}=(1.1\pm0.1)\cdot 10^{-14}\,\mathrm{mol\,Pa^{-1}m^{-1}s^{-1}}$ to calculate the effective permeability coefficient for air $P_{\mathrm{air}}=(5.3\pm0.5)\cdot 10^{-15}\mathrm{mol\,Pa^{-1}m^{-1}s^{-1}}$.
Expressing $\gamma$ in units of moles instead of the volume of air exiting the bubble, we can calculate the thickness of the bubble as
\begin{equation}
	d=\frac{(5.3\pm0.5)\cdot 10^{-15}\,\mathrm{mol\,Pa^{-1}m^{-1}s^{-1}}}{(8.1\pm0.4)\cdot 10^{-8}\,\mathrm{mol\,Pa^{-1}m^{-2}s^{-1}}}= \SI{65\pm9}{nm}.
\end{equation}
This value is consistent with the value calculated from the measured effective refractive index.

\clearpage	




\section{Sensitivity to pressure changes}

A pressure difference between the interior and the exterior of the bubble is given by the Laplace law.
For a small change in the external pressure, the size of the bubble changes such that the pressure difference satisfies Laplace condition.
The pressure in the bubble is changed due to the change in the volume of the bubble.
If we assume that the process is isothermal, the pressure inside the bubble after the change, $p_i'$ is:
\begin{equation}
	\frac{p_i'}{p_i}=\frac{V}{V'},
	\label{isothermal}
\end{equation}
where $p_i$ is the initial pressure in the bubble, $V$ and $V'$ are the initial and final volumes of the bubble, respectively.
Since the bubbles are connected to the capillary, the volume of both the capillary ($V_c$, which remains unchanged) and the bubble ($V_b$) must be included in Eq. \eqref{isothermal}, $V=V_c+V_b$ and $V'=V_c+V_b'$.
The Laplace pressure can be written explicitly as the difference between the pressures on both sides of the film:
\begin{align}
	\label{p_L}
	p_L&=\frac{4\sigma}{R}=p_i-p_0,\\ 
	\label{p_L'}
	p_L'&=\frac{4\sigma}{R'}=p_i'-(p_0+\Delta p),
\end{align}
where $p_L$ and $p_L'$ are Laplace pressures before and after the change, $p_0$ is the initial external pressure, and $\Delta p$ is the change in the external pressure.
Combining  Eq. \eqref{isothermal}, \eqref{p_L} and \eqref{p_L'} we obtain the equation that relates the change in external pressure to the change in bubble size :
\begin{equation}
	\Delta p = \frac{V}{V'}p_L - p_L' + p_0\left(\frac{V}{V'}-1\right). \label{dp}
\end{equation}
In the approximation of small radius change, we get the relation between the small size and pressure change:
\begin{equation}
	\frac{\Delta R}{\Delta p} =  R\left[\frac{3}{1+V_c/V_b}\,\left(p_0+\frac{4\sigma}{R}\right) - \frac{4\sigma}{R}\right]^{-1}
\end{equation}


\begin{figure}[htbp]
	\renewcommand\figurename{Supplementary Figure}
	\centering
	\includegraphics[width=15cm]{FigS9.pdf}
	\caption{Calculated critical volume, and minimal and maximal measurable pressures. (a) Critical additional volume relative to the volume of the bubble, at which the bubble becomes unstable. Above this additional volume the bubble collapses due to the Laplace pressure. For a \SI{2}{mm} bubble, this is at approximately 3000. The red circle represents the bubble size used in the experiments. (b) The minimum measurable pressure difference. (c) Relation between minimum and maximum measurable pressures as $V_c/V_b$ is changed. The red circle in b and c represents the maximum $V_c/V_b$ used in the experiments.
	}
	\label{fig9}
\end{figure}



Maximal measurable pressure depends on the additional volume in the capillary $V_c$ and is calculated by solving Eq. \eqref{dp}. It is the maximal pressure at certain $V_c$ at which the solutions exist. At some critical volume $V_c^{crit}$, there is no pressure change $\Delta p$ for which the solution of Eq. \eqref{dp} would exist. At $V_c$ higher than $V_c^{crit}$ the bubbles are not stable any more. $V_c^{crit}$ is dependent on the surface tension and size of the bubble (Fig. S\ref{fig9}a).


	\clearpage

	
	%\bibliographystyle{unsrt}
	\bibliographystyle{naturemag}
	\bibliography{references}
	
	
\end{document}


